% ---
% id: receiver/rem:exactness-scope
% kind: rem
% title: "Scope of the exactness rule"
% summary: >-
%   The eigenvalue-slack assumption is the one relation between a floating-point
%   and an exact number the receiver asserts; enclosures and rigorous weights
%   are proved. A claim outside the exception of
%   [[thm:exact-and-evaluated]] may still mention a receiver tolerance, in a
%   cross-reference, a remark or an example.
% prominence: standard
% status: accepted
% uses: []
% mentions: [ass:eig, def:pmcex, thm:exact-and-evaluated, thm:enclosed-blocks, thm:enclosed-weight, req:evaluation-label, def:evaluation-model, def:tolmap, def:bounddata, receiver/detector-imperfections/def:darkadm, receiver/detector-imperfections/rem:corrdark, rem:noslack, rem:exacteig, receiver/component-catalogue/ex:bb84exact]
% assumes: []
% refs: []
% provenance: newly-derived
% lean:
% ---

\begin{remark}\label{rem:exactness-scope}
The one relation between a floating-point number and an exact one that this
contract does assert is Assumption~\ref{ass:eig}; it is a separate assumption
because it is invoked only on the opt-in route of Definition~\ref{def:pmcex},
evaluated in floating point.
The enclosures of the photon-number blocks
(Theorem~\ref{thm:enclosed-blocks}) relate rationals computed in exact
arithmetic to the model, and so do the rigorously evaluated weights
(Theorem~\ref{thm:enclosed-weight}); they are proved, not assumed.
An artifact carries the rigorous evaluation label exactly where every weight
it emits was evaluated rigorously, and the floating-point label otherwise; no
path in this module produces another label
(Theorem~\ref{thm:exact-and-evaluated}\,(iii),
Requirement~\ref{req:evaluation-label}).

Item~(i) of Theorem~\ref{thm:exact-and-evaluated} does not say that a claim
outside its exception never \emph{mentions} a receiver tolerance: such claims
do, in cross-references to Definition~\ref{def:tolmap} recording which
declared number decides an exact condition in the evaluation model of
Definition~\ref{def:evaluation-model} (Definition~\ref{def:bounddata},
Definition~\ref{def:darkadm}, Remark~\ref{rem:corrdark}), and in remarks and
an example that quantify the effect of a rule stated in that theorem
(Remark~\ref{rem:noslack}, Remark~\ref{rem:exacteig},
Example~\ref{ex:bb84exact}); a cross-reference asserts nothing. Nor does the
item \emph{enumerate} the exceptions: the rule decides each claim by reading
that claim's own statement.
\end{remark}
